% !TeX program = pdflatex
% !BIB program = biber
% Overleaf: set Compiler = pdfLaTeX, Main document = main.tex
%
% Companion paper to overleaf_pef_article/.
% Solo-authored mathematical infrastructure for the empirical PEF framework:
% canonical form, geometry, partition-function structure, Fisher--Rao coordinates,
% and the corresponding numerical validation.
%
% NOTE: This project is intentionally scaffolded only.  The empirical paper
% (../overleaf_pef_article/) is to be finalised and submitted first; this
% companion will be drafted in parallel and preprinted around the same time.

\documentclass[11pt,a4paper]{article}

\usepackage[british]{babel}
\usepackage[T1]{fontenc}
\usepackage{amsmath,amssymb,amsthm,mathtools}
\usepackage{bbm}
\usepackage{booktabs}
\usepackage{graphicx}
\usepackage{geometry}
\usepackage{microtype}
\usepackage{xcolor}
\usepackage{csquotes}
\usepackage{hyperref}
\usepackage[nameinlink,noabbrev]{cleveref}

\geometry{margin=1in}

\hypersetup{
  colorlinks=true,
  linkcolor=blue!50!black,
  citecolor=blue!50!black,
  urlcolor=blue!50!black
}

% Vancouver-style: numbered citations in order of first appearance.
\usepackage[style=numeric,backend=biber,sorting=none,giveninits=true,maxbibnames=99]{biblatex}
\addbibresource{references.bib}

% --- Notation (UK English in prose; LaTeX maths per OVERLEAF_BEST_PRACTICES) ---
\newcommand{\Var}{\operatorname{Var}}
\newcommand{\Cov}{\operatorname{Cov}}
\DeclareMathOperator{\Corr}{Corr}
\newcommand{\ind}{\mathbbm{1}}
\newcommand{\sech}{\operatorname{sech}}
\newcommand{\arctanh}{\operatorname{arctanh}}
\newcommand{\Real}{\operatorname{Re}}

% Theorem environments (kept light; numbering scheme chosen so the companion
% paper carries its own theorem stream independent of the empirical paper).
\theoremstyle{plain}
\newtheorem{theorem}{Theorem}[section]
\newtheorem{proposition}[theorem]{Proposition}
\newtheorem{lemma}[theorem]{Lemma}
\newtheorem{corollary}[theorem]{Corollary}

\theoremstyle{definition}
\newtheorem{definition}[theorem]{Definition}

\theoremstyle{remark}
\newtheorem{remark}[theorem]{Remark}

\title{The Geometry of Paired Efficiency:\\
       Canonical Form, Information-Geometric Structure,\\
       and Cross-Domain Pooling for Fisher's Generalised Formula}
\author{Rowan Brown}
\date{}

\begin{document}
\maketitle

% TODO: draft last, once Sections 2--7 are in their final form.
% Target length: 200--250 words.  No equations (per OVERLEAF_BEST_PRACTICES).
%
% Required points:
%  - The companion empirical paper (Brown et al.) generalises Fisher's paired
%    efficiency to unequal variances through the Paired Efficiency Factor
%    eta = (1+kappa)/(1+kappa - 2*sqrt(kappa)*rho).
%  - This paper shows that eta carries non-obvious geometric and
%    information-theoretic structure: a canonical form on the half-strip
%    (|tau|, rho) with tau = (1/2) log kappa; an involution kappa <-> 1/kappa;
%    a Moebius linearisation and Chebyshev/Poisson expansion; a Riemann-sphere
%    realisation eta = 1/(1 - cos sigma); a cumulant generating function
%    log eta = -log(1 - rho sech tau) for a one-parameter geometric
%    exponential family; and a Fisher-Rao variance-stabilising coordinate
%    psi = 2 arctanh sqrt(rho sech tau).
%  - These results are distribution-free; the bivariate-normal mutual-information
%    relationship in the empirical companion is the Gaussian specialisation.
%  - Numerical validation across the six datasets of the empirical companion
%    confirms (i) the kappa <-> 1/kappa symmetry holds at the estimator and
%    pipeline level, (ii) the cumulant identity holds within sampling noise
%    away from rho = 0, (iii) psi stabilises cross-domain variance, and
%    (iv) the predicted regime change at rho = 0 is observable in the
%    machine-learning residuals.
%  - Implication: the four-quadrant taxonomy of the empirical paper is a
%    2:1 covering of a substantive two-regime structure indexed by sign(rho)
%    and |tau|; psi is the natural meta-analytic coordinate.

\begin{abstract}
% PLACEHOLDER --- write after Sections 2--7 are stable.
\end{abstract}

\section{Introduction}
\label{sec:intro}

% Intended content (~2 pages):
%
%  1.1 Fisher's paired efficiency and its generalisation.
%      One paragraph recapping Fisher's 1/(1 - rho) and the PEF generalisation
%      eta = (1 + kappa)/(1 + kappa - 2*sqrt(kappa)*rho).  Cite the empirical
%      companion paper for the cross-domain validation; this paper studies
%      the structure of the formula itself.
%
%  1.2 Why the formula has more structure than it appears to.
%      Motivate the half-strip reparameterisation tau = (1/2) log kappa and
%      the canonical form eta = cosh tau / (cosh tau - rho).  Explain that
%      everything in Sections 2--6 is a consequence of this single identity.
%
%  1.3 Contributions.
%      List six contributions:
%        (i)   canonical form on the half-strip and the kappa <-> 1/kappa
%              involution;
%        (ii)  Moebius linearisation 1/eta = 1 - rho sech tau and the
%              Chebyshev/Poisson series;
%        (iii) Riemann-sphere realisation eta = 1/(1 - cos sigma) and the
%              identification of the Fisher regime (kappa = 1) with the
%              sphere's equator;
%        (iv)  partition-function reading log eta = -log(1 - rho sech tau)
%              and the latent geometric exponential family;
%        (v)   Fisher-Rao geometry, the variance-stabilising coordinate psi,
%              and the regime change at rho = 0;
%        (vi)  numerical validation of (i)--(v) on the six datasets of the
%              empirical companion.
%
%  1.4 Relation to the empirical companion paper.
%      Forward reference: empirical companion uses Sections 2--6 results as
%      cited infrastructure; this paper uses Section 7's numerical results to
%      ground the geometry.  Each paper is independently readable.
%
%  1.5 Outline.
%      Pointer to Sections 2--8.

% TODO: drafting.

\section{Canonical Form and the \texorpdfstring{$\kappa \leftrightarrow 1/\kappa$}{kappa <-> 1/kappa} Symmetry}
\label{sec:canonical}

% Intended content (~2 pages):
%
%  2.1 The change of variables.
%      Derive tau = (1/2) log kappa, sqrt(kappa) + 1/sqrt(kappa) = 2 cosh tau,
%      and the canonical form
%        eta(tau, rho) = cosh tau / (cosh tau - rho).
%      State explicitly that everything in Sections 3--6 is a consequence
%      of this identity.
%
%  2.2 The involution.
%      Theorem: eta(kappa, rho) = eta(1/kappa, rho) for all kappa > 0 and
%      rho in (-1, 1) such that the denominator is non-zero.  One-line proof
%      from the canonical form.  Corollary: the substantive moduli space is
%      tau >= 0 (equivalently kappa >= 1), modulo the involution kappa -> 1/kappa.
%
%  2.3 Domain and singularity locus.
%      The denominator vanishes on rho = cosh tau, which intersects the
%      physical domain rho in [-1, 1] only at (tau, rho) = (0, 1).  Local
%      parabolic blow-up
%        eta ~ 1 / ((1 - rho) + (1/2) tau^2).
%      Iso-eta contours near the pole are parabolas, not circles.
%
%  2.4 Differential structure.
%      d_rho eta > 0 everywhere; d_tau eta = 0 on tau = 0 and on rho = 0;
%      sensitivity to correlation is maximised at tau = 0 and decays as
%      |tau| -> infinity.  Mixed derivative analysis.  No interior critical
%      points.
%
%  2.5 Special values.
%      Tabulate boundary values eta(rho = +1, tau), eta(rho = -1, tau),
%      eta(kappa -> 0 or infinity), and the reciprocal identity
%        eta(rho = +1, tau) * eta(rho = -1, tau) = 1.
%      Flag this as a residual echo of the kappa <-> 1/kappa duality.
%
%  2.6 Consequence for the empirical four-quadrant taxonomy.
%      The four-quadrant taxonomy of the empirical companion is a 2:1 covering
%      of the substantive half-strip (|tau|, rho).  Falsifiable: pooling
%      kappa -> max(kappa, 1/kappa) across the six empirical datasets must
%      collapse Q1 <-> Q2 and Q3 <-> Q4 within bootstrap tolerance.  Validated
%      in Section 7.

% TODO: drafting.

\section{M\"obius Linearisation and the Chebyshev/Poisson Expansion}
\label{sec:mobius}

\subsection{M\"obius Linearisation}
\label{sec:mobius_linear}

From the canonical form \cref{eq:canonical}, taking the reciprocal gives immediately:
\begin{equation}
  \frac{1}{\eta} = 1 - \rho\,\sech\tau.
  \label{eq:reciprocal_linear}
\end{equation}
The map $(\tau,\rho) \mapsto 1/\eta$ is therefore \emph{affine} in $\rho$ for fixed $\tau$, and affine in $\sech\tau$ for fixed $\rho$.  Define the rescaled correlation
\begin{equation}
  \tilde\rho = \rho\,\sech\tau = \cos\sigma,
  \qquad\tilde\eta = 1 - 1/\eta.
  \label{eq:mobius_coords}
\end{equation}
In the coordinates $(\tau,\tilde\rho)$, the level set $\{1/\eta = c\}$ becomes $\tilde\rho = 1 - c$: a horizontal plane.  The curvature of iso-$\eta$ contours in the original $(\tau,\rho)$-coordinates is entirely an artefact of viewing the flat surface $\tilde\eta = \tilde\rho$ through the nonlinear coordinate change $\tilde\rho = \rho\,\sech\tau$.

The iso-$\eta$ contours are made explicit by setting $1/\eta = c$ in \cref{eq:reciprocal_linear}:
\begin{equation}
  \rho = (1-c)\cosh\tau, \qquad c \in (0,1)\cup(1,\infty).
  \label{eq:iso_eta}
\end{equation}
For $\eta > 1$ ($c < 1$), contours are cosh curves that open upward from $(0,1-c)$; for $\eta < 1$ ($c > 1$), they open downward from $(0,1-c) < 0$.  At $\eta = 1$ the contour degenerates to $\rho = 0$.

\subsection{Poisson-Kernel Form}

Setting $r = \sqrt{\kappa} = e^{\tau}$ and $\rho = \cos\varphi$ with $\varphi \in (0,\pi)$, the original PEF formula \cref{eq:pef_original} becomes
\begin{equation}
  \eta = \frac{1 + r^{2}}{|1 - r\,e^{i\varphi}|^{2}}.
  \label{eq:poisson_form}
\end{equation}
The classical Poisson kernel for the disc is $P_{r}(\varphi) = (1-r^{2})/|1-r\,e^{i\varphi}|^{2}$, so
\begin{equation}
  \eta = P_{r}(\varphi) + \frac{2r^{2}}{|1-r\,e^{i\varphi}|^{2}}
       = P_{r}(\varphi) + 2r^{2}\,\eta/(1+r^{2}).
  \label{eq:poisson_decomposition}
\end{equation}
Rearranging, $\eta\,(1-2r^{2}/(1+r^{2})) = P_{r}(\varphi)$, i.e.\ $\eta \cdot (1-r^{2})/(1+r^{2}) = P_{r}(\varphi)$.  Equivalently, $\eta = (1+r^{2})/(1-r^{2})\cdot P_{r}(\varphi)$, which gives a simple ratio representation:
\begin{equation}
  \frac{\eta(\tau,\rho)}{P_{r}(\varphi)}
  = \frac{1+r^{2}}{1-r^{2}}
  = \frac{\cosh(2\tau)}{\sinh(2\tau)}\cdot\sinh(2\tau)/(1+e^{-4\tau})
  = \coth(\tau)\big|_{\tau>0}.
  \label{eq:eta_poisson_ratio}
\end{equation}
In particular, $\eta \ge P_{r}(\varphi)$ everywhere with equality only as $\tau \to 0$.

\subsection{Geometric and Chebyshev Expansions}
\label{sec:chebyshev}

From the affine reciprocal \cref{eq:reciprocal_linear}, write $w = \rho\,\sech\tau$.  Whenever $|w| < 1$ (equivalently $|\rho| < \sqrt{\kappa}$ on the fundamental domain $\tau \ge 0$),
\begin{equation}
  \eta(\tau,\rho)
  = \frac{1}{1-w}
  = \sum_{k=0}^{\infty}(\rho\,\sech\tau)^{k}.
  \label{eq:geometric_series}
\end{equation}
Convergence is geometric since $|\rho\,\sech\tau| \le \sech\tau < 1$ for $\tau > 0$ and $|\rho| < 1$.

\begin{proposition}[Geometric series]
\label{prop:geometric_series}
For $|\rho\,\sech\tau| < 1$,
\begin{equation}
  \eta(\tau,\rho) = \sum_{k=0}^{\infty}(\rho\,\sech\tau)^{k},
  \label{eq:chebyshev_exp}
\end{equation}
with uniform convergence on compact subsets of $\{|\rho| < 1,\,\tau > 0\}$.
\end{proposition}

\begin{proof}
Immediate from \cref{eq:reciprocal_linear} and the identity $1/(1-w) = \sum_{k=0}^{\infty}w^{k}$ for $|w| < 1$.
\end{proof}

\begin{remark}[Chebyshev-$U$ cognate]
In the original variables with $r = \sqrt{\kappa}$, \cref{eq:pef_original} gives $\eta = (1+r^{2})/(1-2r\rho+r^{2})$.  The denominator is the generating function for Chebyshev polynomials of the second kind \parencite{rivlin1990}:
\[
  \frac{1}{1-2r\rho+r^{2}} = \sum_{n=0}^{\infty}U_{n}(\rho)\,r^{n},
  \qquad |r| < 1,
\]
so $\eta = (1+r^{2})\sum_{n=0}^{\infty}U_{n}(\rho)\,r^{n}$ on the branch $\kappa < 1$.  On the fundamental domain $\kappa \ge 1$, use \cref{eq:geometric_series} instead.  Both are classical rearrangements of the same rational function; the novelty is their interpretation for paired efficiency, not the series identities themselves.
\end{remark}

\begin{remark}
At $\tau = 0$ (Fisher's case), $\sech\tau = 1$ and \cref{eq:geometric_series} becomes $\eta = 1/(1-\rho) = \sum_{k=0}^{\infty}\rho^{k}$, the geometric series for the Fisher efficiency.
\end{remark}

\subsection{Connection to Harmonic Analysis}
\label{sec:harmonic}

The Poisson-kernel representation \cref{eq:poisson_form} places $\eta$ in the framework of bounded analytic functions on the disc \parencite{garnett2007}.  The function $f(z) = (1+|z|^{2})/|1-z|^{2}$ is harmonic in $z = re^{i\varphi}$ for $r < 1$ (with $r = \sqrt{\kappa}$ playing the role of the radial variable and $\varphi = \arccos\rho$ the angular variable).  The boundary behaviour as $r \to 1$ recovers Fisher's formula $\eta \to 1/(1-\rho)$ on the unit circle $|z| = 1$.

This harmonic-analytic perspective suggests a richer operator-theoretic treatment (Toeplitz operators, Hardy-space norms) as future work, but we do not pursue this here beyond noting that $\eta$ is a positive harmonic function on the disc that extends to a meromorphic function on the Riemann sphere with a single pole at $z = 1$ (corresponding to $(\kappa,\rho) = (1,1)$), confirming the geometry of \cref{sec:sphere}.

\section{$\eta$ as Inverse Spherical Chord}
\label{sec:sphere}

% Intended content (~3 pages --- this is one of the two genuinely novel sections).
%
%  4.1 The lift.
%      Definition: for (kappa, rho) in (0, infinity) x (-1, 1), set
%        z = sqrt(kappa) * exp(i arccos rho)
%      and let pi^{-1}: C -> S^2 be the inverse stereographic projection from
%      the south pole.  This realises every (kappa, rho) as a point on the
%      Riemann sphere.
%
%  4.2 The pole at independence.
%      The point P(1) = (1, 0, 0) on S^2 corresponds to (kappa, rho) = (1, 1):
%      Fisher's perfectly-paired, equal-variance limit, where eta -> infinity.
%      Call P(1) the marked pole.
%
%  4.3 The angular distance identity.
%      Theorem (sphere realisation):
%        eta(kappa, rho) = 1 / (1 - cos sigma) = (1/2) csc^2(sigma/2),
%      where sigma is the angular (spherical) distance from pi^{-1}(z) to
%      the marked pole P(1).  Proof: direct computation, using
%      cos sigma = 2 sqrt(kappa) rho / (1 + kappa).
%
%  4.4 Geometric translations of the empirical taxonomy.
%      - Fisher's classical line kappa = 1 is the equator.
%      - The involution kappa <-> 1/kappa is reflection across the equator.
%      - Iso-eta contours are spherical caps centred on P(1).
%      - The empirical "competitive, negative-rho, high-asymmetry" Q4 regime
%        is literally the antipodal hemisphere to the marked pole.
%
%  4.5 Why this matters for the empirical paper.
%      The four-quadrant taxonomy of the empirical companion is a *coordinate
%      artefact* of viewing a spherical object through a planar (kappa, rho)
%      chart.  The substantive structure is the angular distance sigma to
%      independence; the four quadrants are just the four sign-patterns of
%      (kappa - 1, rho), which on the sphere collapse to a single
%      latitude/longitude coordinate up to the kappa <-> 1/kappa reflection.
%
%  4.6 Caveats.
%      The stereographic lift is not unique (a different choice of base
%      point gives a different chart); we choose the one that makes P(1)
%      the marked pole because it is the only point in the closure of the
%      physical domain where eta diverges.

% TODO: drafting.  Literature scan required: check whether this projective
% realisation of the variance-ratio / correlation pair has appeared in
% statistical-shape-analysis or random-matrix-eigenvalue work.  My current
% belief is that the *specific* realisation with marked pole at (1, 1) is
% novel for this family, but the more general fact that 2x2 correlation
% matrices lift to the sphere is well-known and must be cited.

\section{Partition-Function Structure}
\label{sec:partition}

\subsection{The Cumulant-Generating Identity}

The canonical form \cref{eq:canonical} yields an immediate log-identity:
\begin{equation}
  \log\eta(\tau,\rho)
  = -\log\!\bigl(1 - \rho\,\sech\tau\bigr).
  \label{eq:logeta_identity}
\end{equation}
Setting $u = \rho\,\sech\tau$ and $\theta = \log u$ for $u \in (0,1)$, the right-hand side becomes
\[
  -\log(1 - e^{\theta}) = \sum_{n=1}^{\infty}\frac{e^{n\theta}}{n},
\]
which is the log-partition function (cumulant generating function at zero argument) of the \emph{logarithmic series distribution}. A more elementary identification comes from the geometric series:
\[
  \eta = \frac{1}{1 - u} = \sum_{n=0}^{\infty}u^{n},\qquad u \in (0,1),
\]
recognising $u^n = e^{n\theta}$ as the natural-parameter exponential-family weighting of a latent count $n \in \{0,1,2,\dotsc\}$.

\begin{theorem}[Geometric Exponential Family]
\label{thm:exp_family}
For $\rho\,\sech\tau \in (0,1)$, define the natural parameter $\theta = \log(\rho\,\sech\tau) < 0$ and the log-partition function $A(\theta) = \log\eta$.  Then
\begin{equation}
  p(n;\theta) = (1-e^{\theta})\,e^{n\theta},\qquad n \in \{0,1,2,\dotsc\},
  \label{eq:geo_family}
\end{equation}
defines a one-parameter exponential family with sufficient statistic $N$, base measure $h(n) = 1$, and $A(\theta) = -\log(1-e^{\theta}) = \log\eta$.  The PEF $\eta(\tau,\rho)$ is the partition function of this family, evaluated at the effective parameter $\theta = \log(\rho\,\sech\tau)$ determined by the pair $(\tau,\rho)$.
\end{theorem}

\begin{proof}
Normalisation: $\sum_{n=0}^{\infty}(1-e^{\theta})e^{n\theta} = (1-e^{\theta})/(1-e^{\theta}) = 1$ for $e^{\theta} \in (0,1)$.  The cumulant generating function is $\log\mathbb{E}[e^{tN}] = -\log(1-e^{t+\theta}) + \log(1-e^{\theta})$, so $A(\theta) = -\log(1-e^{\theta}) = \log\eta$ is the log-partition function evaluated at $t=0$.
\end{proof}

\subsection{First and Second Cumulants}
\label{sec:cumulants}

The standard exponential-family relations $\mathbb{E}[N] = A'(\theta)$ and $\Var(N) = A''(\theta)$ give:

\begin{proposition}[Cumulant Identities]
\label{prop:cumulants}
\[
  \mathbb{E}[N] = A'(\theta) = \frac{e^{\theta}}{1-e^{\theta}} = \eta - 1,
  \qquad
  \Var(N) = A''(\theta) = \frac{e^{\theta}}{(1-e^{\theta})^{2}} = \eta(\eta-1).
\]
\end{proposition}

\begin{proof}
$A'(\theta) = e^{\theta}/(1-e^{\theta})$.  Since $e^{\theta} = u = \rho\,\sech\tau$ and $1-u = 1/\eta$, we get $A'(\theta) = u\eta = \eta - 1/\eta \cdot \eta = \eta - 1$.  For the second derivative, $A''(\theta) = e^{\theta}/(1-e^{\theta})^{2} = u\eta^{2} = (\eta-1)\eta$.
\end{proof}

The identities admit a direct empirical reading.  The quantity $\eta - 1$ is the \emph{signal}: for $\eta > 1$ it is positive, measuring the fractional gain from pairing; for $\eta < 1$ it is negative, measuring the fractional loss.  The quantity $\eta(\eta-1)$ is the \emph{curvature of the log-partition function}, which determines how rapidly $A$ changes with the natural parameter: it is the Fisher information of the latent family, and it sets the scale of fluctuations of $\hat{\eta}$ around its expectation.  Both quantities are properties of $\eta$ alone, valid for any $(\kappa,\rho)$ such that $\rho\,\sech\tau \in (0,1)$, with no additional distributional assumptions beyond the PEF definition.

\subsection{Regime Change at \texorpdfstring{$\rho = 0$}{rho = 0}}
\label{sec:regime}

The geometric-family identification requires $\rho\,\sech\tau \in (0,1)$, i.e.\ $\rho > 0$.  For $\rho < 0$, the series $\sum_{n=0}^{\infty}(\rho\,\sech\tau)^{n}$ becomes sign-alternating: even-$n$ terms contribute positively, odd-$n$ terms negatively.  The closed form $\eta = 1/(1-\rho\,\sech\tau)$ remains well-defined and finite (since $|\rho\,\sech\tau| < 1$ in the physical domain), but the partition-function reading breaks down because $\theta = \log(\rho\,\sech\tau)$ is undefined for $\rho < 0$.

There is a symmetry that partially restores the picture: write $\eta(\tau,\rho)$ for $\rho < 0$ as
\[
  \eta(\tau,\rho)
  = \frac{\cosh\tau}{\cosh\tau + |\rho|}
  = 1 - \frac{|\rho|}{\cosh\tau + |\rho|} < 1.
\]
Setting $\tilde{u} = |\rho|\,\sech\tau \in (0,1)$, we have $\eta = 1 - \tilde{u}/(1+\tilde{u}) = 1/(1+\tilde{u})$, which is distinct from the $\rho > 0$ form $1/(1-u)$.  The two branches meet at $\eta = 1$ when $\rho = 0$.  The partition function structure is therefore intrinsically chiral: it requires a sign choice at $\rho = 0$ before the latent-geometric reading applies.

This is the mathematical origin of the regime change that the empirical paper terms the `efficiency--power tension' in competitive settings.  The prediction is concrete: statistics and machine-learning diagnostics that are smooth functions of $\eta$ will show a detectable slope discontinuity at $\rho = 0$, reflecting the change in the curvature $A''(\theta)$ across the boundary.  \Cref{sec:val_regime} tests this empirically.

\subsection{Distribution-Free Scope}
\label{sec:distrib_free}

The Gaussian mutual-information link of the empirical companion \parencite{brownPEFempirical} requires bivariate normality.  The cumulant identities of \cref{sec:cumulants} do not: they follow from the algebraic form of $\eta$ alone.  Specifically, \cref{prop:cumulants} uses no properties of the original data-generating distribution beyond the definition of $\kappa$ and $\rho$ as (population) variance ratio and correlation.  The result is:
\begin{itemize}
  \item the `signal' quantity $\eta-1$ and the `information' quantity $\eta(\eta-1)$ are \emph{distribution-free} properties of the efficiency landscape;
  \item the Gaussian mutual information $I(X;Y) = -\tfrac{1}{2}\log(1-\rho^{2})$ is a $(\kappa,\rho)$-specific instance arising only when $\kappa = 1$ (Fisher's case) under joint normality;
  \item the partition-function form $\log\eta = -\log(1-\rho\,\sech\tau)$ is the natural generalisation that replaces the Gaussian special case with a formula that is valid across all variance ratios and correlations without distributional assumptions.
\end{itemize}

\subsection{The Latent Count}
\label{sec:latent_count}

The random variable $N$ in \cref{thm:exp_family} has no direct interpretation as a counting variable in the original $(X_{\mathrm{A}},X_{\mathrm{B}})$ problem.  It is a \emph{latent} sufficient statistic in the sense of exponential-family theory \parencite{brown1986,coverThomas2006}: a formal object whose existence is implied by the algebraic form of $A(\theta)$, not by any physical counting process.  Analogous latent constructions appear throughout statistics; the Poisson--Binomial relationship and the Pólya-urn representation of the Dirichlet--Multinomial are comparable examples.

The interpretation we adopt is functional rather than generative: $\mathbb{E}[N] = \eta - 1$ quantifies deviation from the pairing-null ($\eta = 1$, $N \equiv 0$), and $\Var(N) = \eta(\eta-1)$ quantifies the sensitivity of the log-partition function to changes in the effective parameter $\theta = \log(\rho\,\sech\tau)$.  These are empirically accessible through the estimator $\hat{\eta}$, and their validity is independent of the data-generating mechanism.

\section{Fisher--Rao Geometry and the \texorpdfstring{$\psi$}{psi} Coordinate}
\label{sec:fisher_rao}

\subsection{The Variance-Stabilising Coordinate}
\label{sec:psi_def}

The Fisher information of the geometric exponential family (\cref{thm:exp_family}) at natural parameter $\theta = \log(\rho\,\sech\tau)$ is
\begin{equation}
  \mathcal{I}(\theta) = A''(\theta) = \eta(\eta-1),
  \label{eq:fisher_info}
\end{equation}
from \cref{prop:cumulants}.  The Fisher--Rao arc-length element on the half-strip is $\mathrm{d}s^{2} = \mathcal{I}(\theta)\,\mathrm{d}\theta^{2}$.  A coordinate $\psi$ is \emph{variance-stabilising} if $\mathrm{d}\psi/\mathrm{d}\theta = \sqrt{\mathcal{I}(\theta)}$, i.e.\ if $\psi$ measures arc length along the Fisher--Rao geodesic from independence.

\begin{proposition}[Variance-Stabilising Coordinate]
\label{prop:psi}
For $\rho\,\sech\tau \in (0,1)$, the coordinate
\begin{equation}
  \psi = 2\,\arctanh\!\sqrt{\rho\,\sech\tau}
       = 2\,\arctanh\!\sqrt{\frac{2\sqrt{\kappa}\,\rho}{1+\kappa}}
  \label{eq:psi_def}
\end{equation}
satisfies $\mathrm{d}\psi = \sqrt{\mathcal{I}(\theta)}\,\mathrm{d}\theta$ and hence gives the Fisher--Rao geodesic distance from the independence point $(\tau,\rho) \in \{\rho = 0\}$.
\end{proposition}

\begin{proof}
Let $u = \rho\,\sech\tau = e^{\theta}$.  Then $\theta = \log u$ and $\mathrm{d}\theta = \mathrm{d}u/u$.  The arc-length integrand is
\[
  \sqrt{\mathcal{I}(\theta)}\,\mathrm{d}\theta
  = \sqrt{\eta(\eta-1)}\cdot\frac{\mathrm{d}u}{u}
  = \frac{\sqrt{u/(1-u)^{2}}}{u}\,\mathrm{d}u
  = \frac{\mathrm{d}u}{\sqrt{u}\,(1-u)}.
\]
Substituting $u = v^{2}$ gives $\mathrm{d}u = 2v\,\mathrm{d}v$ and $\sqrt{u} = v$, so the integrand becomes $2\,\mathrm{d}v/(1-v^{2}) = 2\,\mathrm{d}\!\arctanh(v)$.  Integrating from $u=0$ (independence) gives $\psi = 2\,\arctanh(\sqrt{u}) = 2\,\arctanh\!\sqrt{\rho\,\sech\tau}$.
\end{proof}

\subsection{The Variance-Stabilising Property}
\label{sec:psi_var}

Let $\hat\theta$ be the maximum-likelihood estimator of $\theta$ based on $n$ matched observations.  By standard exponential-family theory, $\sqrt{n}(\hat\theta - \theta)\xrightarrow{d}\mathcal{N}(0,1/\mathcal{I}(\theta))$.  Applying the delta method to $\psi = \psi(\theta)$:
\begin{equation}
  \sqrt{n}(\hat\psi - \psi)
  \xrightarrow{d}
  \mathcal{N}\!\Bigl(0,\Bigl(\tfrac{\mathrm{d}\psi}{\mathrm{d}\theta}\Bigr)^{2}/\mathcal{I}(\theta)\Bigr)
  = \mathcal{N}(0,1),
  \label{eq:psi_clt}
\end{equation}
since $(\mathrm{d}\psi/\mathrm{d}\theta)^{2} = \mathcal{I}(\theta)$ by construction.  Therefore $\Var(\hat\psi) \approx 1/n$ \emph{uniformly} across the half-strip $(\tau,\rho)$, independent of $\eta$.  By contrast, the delta method applied to $\hat\eta$ gives $\Var(\hat\eta) \approx \eta^{2}(\eta-1)/n$, which varies from near zero (when $\eta \approx 1$) to large values (when $\eta \gg 1$).

\begin{remark}
The coordinate $\psi$ is the PEF analogue of Fisher's $z = \arctanh\rho$ \parencite{fisher1921} for the bivariate-normal correlation: both are Fisher--Rao arc lengths \parencite{amari2016} for their respective one-parameter families.  The PEF family has a richer parameter space $(\tau,\rho)$, and $\psi$ maps the two-dimensional family down to a one-dimensional arc length via the projection $u = \rho\,\sech\tau$.
\end{remark}

\subsection{Regime Change at \texorpdfstring{$\rho = 0$}{rho = 0}}
\label{sec:psi_regime}

The coordinate $\psi$ is defined for $u = \rho\,\sech\tau \in (0,1)$, i.e.\ $\rho > 0$.  For $\rho < 0$ the quantity $\rho\,\sech\tau$ is negative, and $\arctanh$ of a negative square root is not real.  The natural extension is to define
\begin{equation}
  \psi = \operatorname{sign}(\rho)\cdot 2\,\arctanh\!\sqrt{|\rho|\,\sech\tau},
  \label{eq:psi_signed}
\end{equation}
so that $\psi \in \mathbb{R}$ for all $(\tau,\rho)$ in the physical domain and $\psi = 0 \Leftrightarrow \rho = 0$.

The change of sign at $\rho = 0$ is not a coordinate singularity: $\psi$ is continuous through $\rho = 0$.  However, the Fisher information $\mathcal{I}(\theta) = \eta(\eta-1)$ changes sign there as well (from positive for $\eta > 1$ to negative for $\eta < 1$), reflecting the geometric fact that the partition-function reading of \cref{sec:partition} switches branches.  The asymptotic variance $\Var(\hat\psi) \approx 1/n$ holds in each regime separately; the regime boundary at $\rho = 0$ is a locus where the local curvature of the exponential family changes character.  This is the content of the regime-change LRT test in \cref{sec:val_regime}, implemented following Davies \parencite{davies1987}.

\subsection{Implication for Cross-Domain Meta-Analysis}
\label{sec:psi_meta}

In the six datasets of \textcite{brownPEFempirical}, mean $\eta$ varies from $0.88$ (football) to $3.13$ (finance) --- a factor of 3.6.  On the $\psi$ scale the values are (from \texttt{psi\_per\_domain.csv}, omitting the single healthcare observation): football $\hat\psi = -0.730$, rugby $+0.217$, manufacturing $+0.017$, clinical genomics $+1.869$, finance $+2.251$.  Bootstrap 95\% confidence intervals for $\psi$ have widths of $0.35$--$0.81$ (the wider end for the smaller rugby dataset, $n=283$), consistent with the $1/\sqrt{n}$ scaling of \cref{eq:psi_clt}.

For meta-analytic pooling, the variance-stabilising property implies that a fixed-effects or random-effects model on the $\psi$ scale \parencite{dersimonianLaird1986} will not be dominated by high-$\eta$ domains (whose $\hat\eta$ has large sampling variance) or compressed by low-$\eta$ domains.  Across the positive-$\rho$ domains (finance, clinical genomics, healthcare) the between-domain coefficient of variation of $\psi$ ($\mathrm{CV}_\psi = 0.41$) is close to that of $\eta$ ($\mathrm{CV}_\eta = 0.33$), confirming that $\psi$ resolves between-domain differences to a similar degree as $\eta$ whilst providing more uniform within-domain confidence intervals.

The practical recipe is: (1) estimate $(\hat\kappa, \hat\rho)$ per KPI; (2) compute $\hat\tau = \tfrac{1}{2}\log\hat\kappa$ and $\hat\psi$ via \cref{eq:psi_signed}; (3) bootstrap $\hat\psi$ confidence intervals; (4) pool on the $\psi$ scale using standard meta-analytic procedures; (5) back-transform to the $\eta$ scale for reporting via $\eta = 1/(1 - \tanh^{2}(\psi/2))$ for $\psi > 0$.

\subsection{Comparison with Related Coordinates}
\label{sec:psi_comparison}

Three coordinates serve analogous roles in related problems:
\begin{itemize}
  \item \textbf{Fisher's $z = \arctanh\rho$ \parencite{fisher1921}.}  Variance-stabilising for the bivariate-normal correlation at equal variances ($\kappa = 1$).  Small-sample bias corrections for $\hat\rho$ are classical \parencite{olkinPratt1958}; an analogous correction for $\hat\psi$ remains open.  Recovers from $\psi$ at $\tau = 0$: $\psi = 2\,\arctanh\!\sqrt{\rho} \ne \arctanh\rho$ in general, but at $\rho \to 0$, $\psi \approx 2\sqrt{\rho} \approx 2z/\sqrt{z}$ --- the two are distinct transformations of the same parameter.
  \item \textbf{Logit $\log(p/(1-p))$.}  Variance-stabilising for binomial proportions.  The cumulant structure is analogous to the geometric family here, with $\eta$ playing the role of $1/q$.
  \item \textbf{Sphere angle $\sigma = \arccos(\rho\,\sech\tau)$.}  Measures colatitude from the marked pole on the unit sphere (\cref{sec:sphere}); related to $\psi$ by $\sigma = \arccos(\tanh^{2}(\psi/2))$, equivalently $\psi = 2\,\arctanh\!\sqrt{\cos\sigma}$, the classical inverse once the sphere lift is fixed.
\end{itemize}
Among these, $\psi$ is the natural one for the geometric exponential family identified in \cref{sec:partition}.

\begin{remark}[Non-parametric extension]
Substituting Spearman or Kendall rank correlation for Pearson $\rho$ in \cref{eq:psi_def} gives a distribution-free version of $\psi$ with the same functional form.  The cumulant identities of \cref{prop:cumulants} no longer have a parametric Fisher-information interpretation, but the coordinate itself remains well-defined and the variance-stabilising property holds approximately for large $n$ by the delta method applied to the rank-correlation estimator.  We leave a formal treatment to future work.
\end{remark}

\section{Numerical Validation}
\label{sec:numerical}

% Intended content (~3 pages --- the empirical engine of the paper).
%
%  7.1 Datasets.
%      The six datasets of the empirical companion paper:
%        - URC rugby union (n = 240 games);
%        - English Championship football (n = 552 games);
%        - NHANES 2017 oscillometric blood pressure (N = 10,352);
%        - GSE47462 paired transcriptomics (genes in ranked subset);
%        - S&P 100 / S&P 500 finance (2020--2023);
%        - Bosch OpenML 752 CNC manufacturing telemetry.
%      Cite empirical companion paper for full data provenance.
%
%  7.2 Symmetry tests (three levels).
%      (1) Algebraic verification on a (kappa, rho) grid; expect residual at
%          machine precision (~1e-14).
%      (2) Estimator-level: for each KPI, swap (X, Y) -> (Y, X), recompute
%          eta-hat, verify exact agreement at machine precision.
%      (2.5) Bootstrap joint-distribution test: for each KPI, bootstrap
%          (|tau-hat|, rho-hat) under both label orderings; energy-distance
%          or 2D KS test for equality.  Catches pipeline asymmetries that the
%          point estimate does not.
%      (3) Pipeline-level (deferred from companion paper to here): rerun the
%          entire ML pipeline with (X, Y) swapped on each dataset; verify
%          ML improvement, AUC, KDE-derived statistics are all invariant.
%          Any non-invariance is reported and attributed to a specific
%          implementation choice (KDE bandwidth, standardisation, ddof, etc.).
%
%  7.3 Cumulant identity check.
%      For each KPI, bootstrap eta-hat and the implied theta-hat; finite-
%      difference the derivatives d_theta log eta and d_theta^2 log eta and
%      compare to eta - 1 and eta(eta - 1) respectively.  Localised failures
%      identify regions of (tau, rho) where the latent geometric model fits
%      poorly; this is the *map* of model fit across the empirical landscape.
%
%  7.4 Sphere-distance equivalence.
%      Compute sigma-hat = arccos(2 sqrt(kappa-hat) rho-hat / (1 + kappa-hat))
%      for each KPI and verify eta-hat = 1 / (1 - cos sigma-hat) numerically.
%      Bootstrap confidence bands.  This is the cleanest empirical witness
%      to the projective interpretation and produces the sphere figure
%      (Figure 4 of this paper).
%
%  7.5 psi-stabilisation.
%      For each domain, bootstrap eta-hat and psi-hat at varying sub-sample
%      sizes n.  Verify Var(psi-hat) is approximately constant across
%      (tau, rho) while Var(eta-hat) scales with eta^2 (eta - 1) / n.  This
%      empirically validates psi as the meta-analytic scale.
%
%  7.6 Regime-change LRT at rho = 0.
%      Fit a continuous model and a model with a slope discontinuity at
%      rho = 0 to (ML gain) on (rho, |tau|) across pooled domains; perform
%      a likelihood-ratio test.  The mathematical regime change of
%      Section 5.3 predicts the discontinuous model wins.  Report p-values,
%      effect sizes, and quadrant-level breakdown.
%
%  7.7 Falsifiability summary.
%      Tabulate what each test would falsify if it failed, and what (if
%      anything) failed in the actual run.  This is the section that gives
%      the paper its empirical bite.

% TODO: drafting.  All computations to be run by the corresponding scripts
% in ../overleaf_pef_article/scripts/paper_pipeline/ once those scripts have
% been added (see the empirical paper's plan file).  The companion paper
% imports the resulting CSVs and reports the numbers.

\section{Discussion}
\label{sec:discussion}

% Intended content (~2 pages).
%
%  8.1 What the geometry buys.
%      Summarise: a coordinate-free description of paired efficiency in which
%      the four-quadrant taxonomy of the empirical companion is a coordinate
%      artefact; a distribution-free signal/information pair via the
%      partition-function reading; a meta-analytic scale (psi) on which
%      cross-domain effect sizes are properly comparable; a single
%      Riemann-sphere figure that organises the framework.
%
%  8.2 Information geometry context.
%      The paired-efficiency family sits inside the broader information-
%      geometric program initiated by Rao (1945), Amari (1985), and others.
%      Brief comparison with Fisher information metrics on more familiar
%      exponential families.  This is the section that earns the paper
%      its place in the information-geometry literature.
%
%  8.3 Relation to the empirical companion paper.
%      Concretely: which results in this paper are *used* by the empirical
%      companion (the kappa <-> 1/kappa symmetry, the psi coordinate, the
%      partition-function remark, the regime-change observation), and which
%      are stand-alone contributions of this paper (the canonical form,
%      Moebius linearisation, Chebyshev/Poisson expansion, sphere lift,
%      Fisher--Rao identification).
%
%  8.4 Limitations.
%      - All results are univariate-pair: an extension to multivariate paired
%        efficiency (PEF matrices) is outlined in Section 8.5 but not
%        developed here.
%      - The partition-function reading requires rho sech tau in (0, 1);
%        the rho < 0 regime is mathematically distinct and is treated by
%        sign-carry, not by analytic continuation.
%      - Numerical validation depends on bootstrap inference; finite-sample
%        coverage of the psi-scale intervals is checked but not theoretically
%        derived to high order.
%
%  8.5 Future directions.
%      Six concrete extensions:
%        (i)   multivariate PEF: a matrix-valued generalisation with its own
%              spectral / spherical structure;
%        (ii)  non-parametric psi: substitute rank correlation for rho-hat
%              and derive the corresponding stabilising transform;
%        (iii) operator-theoretic / Hardy-space treatment building on the
%              Poisson-kernel decomposition of Section 3.4;
%        (iv)  small-sample bias correction for eta-hat and psi-hat;
%        (v)   regime-change asymptotics: refined treatment of the rho = 0
%              singularity in the latent geometric family;
%        (vi)  causal-inference connections: when does correlation-based
%              variance reduction preserve causal estimates?
%
%  8.6 Closing.
%      One-paragraph summary: Fisher's generalised paired efficiency formula
%      is, beneath the algebra, a piece of spherical and information geometry
%      with a distribution-free signal/information pair.  The empirical
%      companion paper validates the formula across six domains; this paper
%      explains *why* the validation works as well as it does, and provides
%      the meta-analytic infrastructure for future cross-domain pooling.

% TODO: drafting.


\appendix
\section{Mathematical Appendix}
\label{sec:appendix}

% Intended content (~2 pages of detailed algebra that would clutter the main text).
%
%  A.1 Derivation of the canonical form (full algebra).
%      Step-by-step from eta = (1 + kappa)/(1 + kappa - 2 sqrt(kappa) rho)
%      to eta = cosh tau / (cosh tau - rho).
%
%  A.2 Proof of the kappa <-> 1/kappa involution.
%      Direct substitution.
%
%  A.3 Proof of the cumulant identity log eta = -log(1 - rho sech tau).
%      Differentiate the canonical form; verify the first two cumulants
%      match eta - 1 and eta (eta - 1).
%
%  A.4 Proof of the sphere realisation eta = 1/(1 - cos sigma).
%      Compute cos sigma = 2 sqrt(kappa) rho / (1 + kappa) directly from the
%      stereographic lift.
%
%  A.5 Variance-stabilising property of psi.
%      Delta-method calculation.
%
%  A.6 Chebyshev/Poisson expansion details.
%      Verification of the identity U_n + U_{n-2} = 2 rho U_{n-1} and
%      collapse of the eta expansion.
%
%  A.7 Bootstrap procedures used in Section 7.
%      Pseudocode for the symmetry test (levels 1, 2, 2.5, 3), cumulant
%      identity check, psi-stabilisation, and regime-change LRT.

% TODO: drafting.  These appendix sections are intentionally separated from
% the main text so that the main text remains readable end-to-end without
% appendix detours.


\printbibliography[title={References}]

\end{document}
